\documentclass{chsmc}
\chsmcsetup{
  year = 2017,
  part = 1,
  type = B,
  show-answers = false,
}
\begin{document}

\section{填空题：本大题共 8 小题，每小题 8 分，共 64 分.}

\begin{enumerate}
  \item 在等比数列 $\{a_n\}$ 中，$a_2 = \sqrt{2}$，$a_3 = \sqrt[3]{3}$，
    则 $\dfrac{a_1 + a_{2011}}{a_7 + a_{2017}}$ 的值为 \blankline
  \item 设复数 $z$ 满足 $z + 9 = 10\bar{z} + 22\ii$，则 $|z|$ 的值为 \blankline
  \item 设 $f(x)$ 是定义在 $\mathbf{R}$ 上的函数，若 $f(x) + x^2$ 是奇函数，
    $f(x) + 2^x$ 是偶函数，则 $f(1)$ 的值为 \blankline
  \item 在 $\triangle ABC$ 中，若 $\sin A = 2\sin C$，且三条边 $a$, $b$, $c$
    成等比数列，则 $\cos A$ 的值为 \blankline
  \item 在正四面体 $ABCD$ 中，$E$, $F$ 分别在棱 $AB$, $AC$ 上，
    满足 $BE = 3$，$EF = 4$，且 $EF$ 与面 $BCD$ 平行，
    则 $\triangle DEF$ 的面积为 \blankline
  \item 在平面直角坐标系 $xOy$ 中，点集
    $K = \{(x,y) \mid x, y = -1, 0, 1\}$. 在 $K$ 中随机取出三个点，
    则这三点两两之间距离均不超过 $2$ 的概率为 \blankline
  \item 设 $a$ 为非零实数，在平面直角坐标系 $xOy$ 中，二次曲线
    $x^2 + ay^2 + a^2 = 0$ 的焦距为 $4$，则 $a$ 的值为 \blankline
  \item 若正整数 $a$, $b$, $c$ 满足 $2017 \geqslant 10a \geqslant 100b \geqslant 1000c$，
    则数组 $(a,b,c)$ 的个数为 \blankline
\end{enumerate}

\section{解答题：本大题共 3 个小题，满分 56 分. 解答应写出文字说明、证明过程或演算步骤.}

\begin{enumerate}[resume]
  \item (本题满分 16 分) 设不等式 $|2^x - a| < |5 - 2^x|$ 对所有
    $x \in [1,2]$ 成立，求实数 $a$ 的取值范围.
  \item (本题满分 20 分) 设数列 $\{a_n\}$ 是等差数列，数列 $\{b_n\}$ 满足
    $b_n = a_{n+1}a_{n+2} - a_n^2$，$n = 1, 2, \cdots$.
    \begin{enumerate}[(1)]
      \item 证明：数列 $\{b_n\}$ 也是等差数列；
      \item 设数列 $\{a_n\}$, $\{b_n\}$ 的公差均是 $d \neq 0$，
        并且存在正整数 $s$, $t$，使得 $a_s + b_t$ 是整数，
        求 $|a_1|$ 的最小值.
    \end{enumerate}
  \item (本题满分 20 分) 在平面直角坐标系 $xOy$ 中，曲线
    $C_1\colon y^2 = 4x$，曲线 $C_2\colon (x-4)^2 + y^2 = 8$.
    经过 $C_1$ 上一点 $P$ 作一条倾斜角为 $45°$ 的直线 $l$，
    与 $C_2$ 交于两个不同的点 $Q$, $R$，求 $|PQ|\cdot|PR|$ 的取值范围.
\end{enumerate}

\end{document}